\section{信息瓶颈在压缩与预测之间权衡}
\label{sec:13-Machine-Learning/11-Information-Theoretic-Learning/04}

图像里有对象的形状、背景的纹理、相机的噪声、拍摄的时间戳。如果任务只是识别对象类别，一个好的表示 \(Z\) 似乎应该忘掉与标签无关的细节，同时保留预测 \(Y\) 所需的那部分信息。

但"应该忘掉"这句话太模糊。忘掉多少算够？哪些细节算"无关"？拍这张照片时是上午还是下午——这个信息和类别有关系吗？训练数据里上午拍的都是猫、下午拍的都是狗，那么时间戳就变成了一个强预测特征。可你知道这只是收集数据的 bias，不是猫和狗的真实区别。

把表示学习写成压缩问题：让 \(Z\) 尽可能少地记住 \(X\)，同时尽可能多地保留关于 \(Y\) 的信息。两个目标天生冲突——记住太少就预测不准，记住太多就带着噪声和虚假相关。

信息瓶颈用互信息把这两个愿望精确化：

\[
\min_{p(z\mid x)}\; I(X;Z) - \beta\,I(Z;Y),\qquad \beta\ge 0,
\]

并要求 \(Z-X-Y\) 构成 Markov 链：表示 \(Z\) 只能从输入 \(X\) 生成，给定输入后不得直接读取标签 \(Y\)。\(\beta\) 是权衡旋钮。\(\beta\) 很大时，模型看重预测，愿意让 \(Z\) 记住更多输入信息。\(\beta\) 很小时，模型看重压缩，宁愿牺牲一些预测精度来换取更紧凑的表示。

两种极端先看清。\(\beta\to 0\)：只管压缩，最优解是常数表示——\(I(X;Z)=0\)，但 \(I(Z;Y)=0\)，什么也预测不了。\(\beta\to\infty\)：只管预测，最优解是 \(X\) 的充分统计量——保留关于 \(Y\) 的全部信息，但 \(I(X;Z)\) 可能等于 \(I(X;Y)\) 加上一堆和 \(Y\) 条件独立的噪声。

关键在于中间地带。信息瓶颈寻找的是那条 \(I(X;Z)\)-\(I(Z;Y)\) 曲线上的最优 trade-off：对于给定的压缩水平，能保留多少任务信息；或者反过来，对于给定的任务信息要求，最少需要从输入里记住多少。

信息瓶颈压缩的是给定联合分布下 \(X\) 与 \(Z\) 的统计依赖，不是自动删除"无意义的细节"。在训练分布里，如果背景纹理碰巧与标签相关，保留背景确实会提高 \(I(Z;Y)\)——目标函数鼓励这种行为。从 IB 目标到语义、鲁棒性或泛化，中间需要额外的假设来桥接。

\subsection{数据处理不等式给出了不可逾越的上限}

由互信息的 chain rule，

\[
I(X,Z;Y) = I(X;Y) + I(Z;Y\mid X)
\]

同时也等于

\[
I(Z;Y) + I(X;Y\mid Z).
\]

Markov 条件 \(Z-X-Y\) 使 \(I(Z;Y\mid X)=0\)——给定输入后，表示不能从标签学到额外的信息。因此

\[
I(Z;Y) = I(X;Y) - I(X;Y\mid Z) \le I(X;Y).
\]

对输入作任何处理都不能创造关于标签的新信息。等号成立意味着给定 \(Z\) 后 \(X\) 不再为 \(Y\) 提供额外信息——\(Z\) 是对任务充分的最小表示，它保留了 \(X\) 中关于 \(Y\) 的每一比特。

但这个定理没有说"压缩必然改善泛化"。训练分布中背景若与标签相关，保留背景确实会提高 \(I(Z;Y)\)；部署时关联消失，原目标不会自动预见这种 shift。鲁棒性需要关于环境、因果结构或不变性的额外假设——IB 提供了一个表述框架，但自身不提供这些假设。

\subsection{深度模型中互信息并不容易算}

高维连续变量的互信息需要未知密度——\(I(X;Z)\) 涉及 \(p(x)\)、\(p(z)\)、\(p(x,z)\) 三个未知分布。更棘手的是：若 \(Z=f(X)\) 是确定性连续映射，在常见的理想化条件下 \(I(X;Z)\) 可能发散或者出现与量化精度强烈耦合的病态行为。表示维数下降也不等于互信息按直觉减少——一个 2 维表示可能携带比 100 维表示更多的输入信息，如果那 100 维是高度冗余的。

变分信息瓶颈通常使用随机 encoder \(q_\phi(z\mid x)\) 和简单 prior \(r(z)\)。记聚合表示分布 \(q_\phi(z) = \int q_\phi(z\mid x)p(x)\,\mathrm dx\)，rate 项可写成

\[
I(X;Z) = \E_x\KL(q_\phi(z\mid x)\Vert q_\phi(z)).
\]

聚合后验 \(q_\phi(z)\) 本身难以计算，变分方法用一个固定 prior \(r(z)\)（通常取标准正态）来近似：

\[
\E_x\KL(q_\phi(z\mid x)\Vert r(z)) = I(X;Z) + \KL(q_\phi(z)\Vert r(z)) \ge I(X;Z).
\]

因此 \(\E_x\KL(q_\phi(z\mid x)\Vert r(z))\) 给出了 rate 项 \(I(X;Z)\) 的上界。Decoder 的 log-likelihood \(\E_{q_\phi(z\mid x)}[\log p_\theta(y\mid z)]\) 则给出任务信息 \(I(Z;Y)\) 的下界（至多相差与模型无关的常数）。实际优化的是这些 bounds 的组合，不一定是精确的 IB 目标。

Bound 的松紧、prior 的选择和 Monte Carlo 估计的方差都会改变最终结果。在小型离散问题上可以枚举 \(p(x,y,z)\) 精确计算互信息，用它核对估计器和 bound 是否有系统性偏差。高维 neural estimator 可能有显著的 bias 与 variance——单次曲线平滑不代表数值可信。

一份可靠的实验应明确：
\begin{itemize}
\tightlist
\item
  encoder 是否随机，噪声如何注入；
\item
  哪一项是 \(I(X;Z)\) 的上界、哪一项是 \(I(Z;Y)\) 的下界；
\item
  \(\beta\) 改变的是哪一个近似 rate--relevance 曲线；
\item
  是否出现常数表示等 collapse 退化；
\item
  预测、校准、分布外鲁棒性和真正估计的 rate 是否分别报告。
\end{itemize}

\subsection{信息瓶颈是一种问题表述，不是万能解释}

IB 最有价值的地方是迫使你把表示学习写成两个冲突目标之间的权衡：保留任务相关信息，同时限制输入细节进入表示。这一语言连接了有损压缩、充分统计量、隐私保护和多视图学习——不同领域看似无关的技术，在 IB 框架下共享同一个数学结构。

但"压缩帮助泛化""压缩产生语义""压缩增强鲁棒性"这些论断，从 IB 目标到它们之间每一步都需要额外的条件和实证检验。IB 把"表示应该忘记无关信息"这句话翻译成了可优化的互信息目标，但它自身不判断哪些相关性是虚假的、哪些环境变化需要警惕。

\subsection{一个离散例子揭示三种表示}

令离散输入 \(X=(Y,N)\)，且 \(N\) 与 \(Y\) 独立。第一种表示保留全部 \(X\)：\(I(X;Z)\) 最大，\(I(Z;Y)\) 达到上限 \(I(X;Y)\)，但携带了与任务无关的噪声 \(N\)。第二种表示只保留 \(Y\)：以更小的 rate 保留全部任务信息，丢弃了 \(N\)。第三种是常数表示：\(I(X;Z)=I(Z;Y)=0\)，两端都归零。

现在改动数据生成过程：训练时 \(N\) 与 \(Y\) 相关（比如：收集数据时晴天拍猫、阴天拍狗，\(N\)="天气"碰巧预测 \(Y\)），测试时独立。此时保留 \(N\) 在训练分布上确实会提高 \(I(Z;Y)\)，但在部署时损害表现。精确 IB 目标本身也受训练联合分布限制——它优化的是训练分布上的 rate 和 relevance。变分上界还会在这层限制之外再引入一层近似差距。

更完整的信息论基础见\hyperref[sec:07-Information-Theory/06-Information-and-Learning/03]{``信息瓶颈与任务相关表示''}。
